% Formalization support lemmas.
%
% NOT book material. These declarations were invented in this workspace to
% carry a Lean formalization and have no counterpart in the source text --
% they carry no \dcref. They live here, outside the book chapters, so the
% book blueprint stays timeless mathematics while their \leanok marks and
% Lean anchors are preserved.
%
% Do not add book mathematics to this file, and do not add formalization
% notes to the book chapters.

\begin{lemma}[Parabolic scaling of heat balls]
\label{lem:heat-ball-parabolic-scaling}
\group{Heat Equation}
\uses{def:heat-ball}
\lean{EvansLib.heatParabolicDilation_mem_heatBall_iff,
  EvansLib.heatParabolicDilation_bijective,
  EvansLib.heatParabolicDilation_image_heatBall,
  EvansLib.heatMeanValueWeight_parabolicDilation}
\leanok
Assume $n \geq 1$ and $r>0$. Under the parabolic dilation
\[
(y,s) \longmapsto (x+ry,t+r^2s),
\]
the unit heat ball $E(0,0;1)$ is mapped bijectively onto $E(x,t;r)$, and for
every $(y,s) \in E(0,0;1)$ (so $s<0$) the Watson weight obeys
\[
\frac{|x-(x+ry)|^2}{(t-(t+r^2s))^2}
 = r^{-2}\frac{|y|^2}{s^2}.
\]
\end{lemma}
\begin{proof}
\leanok
For the positive elapsed time $\tau=-s$ and $z=-y$, the Gaussian satisfies the
scaling identity $\Phi(rz,r^2\tau)=r^{-n}\Phi(z,\tau)$. Substitution in the defining
inequalities therefore gives the membership equivalence; the explicit inverse
dilation proves the set equality, and the bijectivity declaration supplies the
corresponding one-to-one correspondence. The weight identity follows by
cancelling the factors $r^2$ in the numerator and denominator.
\end{proof}

\begin{lemma}[Bounded-data heat convolution: integrability, PDE, and initial trace]
\label{lem:heat-convolution-bounded-initial-trace}
\group{Heat Equation}
\uses{def:heat-convolution-solution,lem:integral-fundamental-solution-heat,
  lem:heat-kernel-solves-heat-equation}
\lean{EvansLib.heatSolution_integrable_of_bounded,
  EvansLib.heatSolution_isClassicalAt_of_bounded,
  EvansLib.heatSolution_tendsto_initial_of_bounded,
  EvansLib.heatSolution_bounded_cauchy_package}
\leanok
If $g \in C(\R^n)$ is globally bounded, then the heat-convolution integrand is
integrable for every $t>0$, and the convolution satisfies
$u_t=\Delta u$ on $\R^n\times(0,\infty)$, with the time derivative and every
pure second spatial derivative existing pointwise. Moreover, it attains the initial data:
for every $x^0 \in \R^n$,
\[
\lim_{(x,t)\to(x^0,0^+)} u(x,t)=g(x^0).
\]
The three bounded-data clauses are also bundled in
\texttt{EvansLib.heatSolution\_bounded\_cauchy\_package}; the all-orders
$C^\infty$ assertion for merely bounded data is a separate regularity question.
\end{lemma}
\begin{proof}
\leanok
Choose $M\geq0$ with $|g(y)|\leq M$. Since $\Phi(\cdot,t)$ is integrable for
$t>0$, so is $y\mapsto\Phi(x-y,t)g(y)$ by domination with
$M\Phi(x-y,t)$.

For $t$ in a compact subinterval of $(0,\infty)$, the time derivative of the
kernel is dominated by an integrable function of the form
$C(1+|x-y|^2)e^{-c|x-y|^2}$. Along each spatial coordinate line, the first and
second kernel derivatives admit the same type of locally uniform dominator.
Dominated differentiation therefore moves $\partial_t$ and every
$\partial_{x_j}^2$ under the integral. The pointwise identity
$\Phi_t=\sum_j\partial_{x_j}^2\Phi$ then gives $u_t=\Delta u$.

Fix $x^0$ and $\varepsilon>0$. Continuity at $x^0$ gives $\delta>0$ such that
$|g(y)-g(x^0)|\leq\varepsilon/2$ whenever $|y-x^0|<\delta$. Using
$\int\Phi(z,t)\,dz=1$ and the substitution $z=x-y$, we have
\[
 u(x,t)-g(x^0)=\int_{\R^n}\Phi(z,t)
       \bigl(g(x-z)-g(x^0)\bigr)\,dz.
\]
If $|x-x^0|<\delta/2$, then $|z|<\delta/2$ implies
$|(x-z)-x^0|<\delta$. Splitting the last integral at $|z|=\delta/2$ therefore
gives
\[
 |u(x,t)-g(x^0)|
 \leq \frac{\varepsilon}{2}
   +2M\int_{\{|z|\geq\delta/2\}}\Phi(z,t)\,dz.
\]
The weighted Gaussian tail term on the right tends to zero as $t\downarrow0$;
choose $t>0$ small enough that
$2M\int_{\{|z|\geq\delta/2\}}\Phi(z,t)\,dz<\varepsilon/2$. This proves the asserted
joint limit as $(x,t)\to(x^0,0^+)$.
\end{proof}

\begin{lemma}[Sourced weak radial ODE identity]
\label{lem:sourced-weak-radial-ode-identity}
\group{Harmonic Functions}
\uses{lem:weak-green-identity-compact-test}
\lean{EvansLib.integral_radius_unitSphereRadialIntegral_radialODE_eq_source}
\leanok
Assume $n\geq 1$. Let $u\in C^2(U)$, where $U\subset\R^n$ is open and
$\overline{B(0,R)}\subset U$. For a smooth profile $g$ that vanishes when
$\rho>R^2$, write
\[
S_v(r)=\int_{\partial B(0,1)}v(r\omega)\,dS(\omega).
\]
Then
\[
\int_0^\infty r^{n-1}S_u(r)
  \left(4r^2g''(r^2)+2n g'(r^2)\right)\,dr
=\int_0^\infty r^{n-1}S_{\Delta u}(r)g(r^2)\,dr.
\]
\end{lemma}
\begin{proof}
\leanok
Set $w(z)=g(|z|^2)$. Its support is contained in
$\overline{B(0,R)}$, and
\[
\Delta w(z)=4|z|^2g''(|z|^2)+2n g'(|z|^2)
\]
away from the origin. Green's identity gives
$\int u\,\Delta w=\int (\Delta u)w$. Applying polar coordinates to both
sides yields the asserted identity.
\end{proof}

\begin{lemma}[Anisotropic heat uniqueness on bounded cylinders]
\label{lem:uniqueness-heat-bounded-domain-anisotropic}
\group{Maximum Principles}
\uses{lem:weak-maximum-principle-heat-anisotropic}
\lean{EvansLib.eqOn_closure_of_eqOn_parabolicBoundary_anisotropic}
\leanok
Let $U \subset \R^n$ be open, bounded and nonempty, and let $T>0$. If two
continuous functions on $\overline{U_T}$ have one time derivative and two pure
spatial derivatives on every coordinate line at interior times, solve the same
heat equation with source $f$ there, and agree on the parabolic boundary, then
they agree on $\overline{U_T}$. This is the bounded-domain uniqueness argument
under Evans's $C^2_1$-style line regularity: subtract the solutions, observe
that the source cancels, and apply the anisotropic weak maximum principle to the
difference and its negative.
\end{lemma}
\begin{proof}
\leanok
The line derivatives are stable under subtraction, and the common source cancels
in the line-form heat equations. The anisotropic weak maximum principle bounds
the difference above by its parabolic-boundary value, while the same argument
for the negative difference bounds it below. Boundary agreement therefore gives
equality throughout the closed cylinder.
\end{proof}

\begin{lemma}[One-dimensional Cauchy uniqueness for waves]
\label{lem:wave-cauchy-uniqueness-1d}
\group{Energy Methods}
\lean{EvansLib.wave_cauchy_unique_1d}
\leanok
Let $u,v \in C^2(\R^2)$ solve the one-dimensional wave equation on all of
$\R^2$. If $u=v$ and $u_t=v_t$ on $\R\times\{t=0\}$, then $u=v$ everywhere.
\end{lemma}
\begin{proof}
\leanok
Subtract the two solutions and write $w=u-v$. The resulting wave has zero
displacement and velocity at time zero. Differentiating the identity
$w(x,0)=0$ in $x$ gives $w_x(x,0)=0$, so both characteristic Riemann
invariants $w_t\pm w_x$ vanish on the initial line. The wave equation makes
these invariants constant along their characteristic lines, hence both vanish
everywhere. Thus $w_t=0$; constancy along vertical time lines and $w(x,0)=0$
give $w\equiv0$.
\end{proof}

\begin{lemma}[One-dimensional finite propagation]
\label{lem:wave-finite-propagation-1d}
\group{Energy Methods}
\lean{EvansLib.wave_finite_propagation_1d}
\leanok
Let $u \in C^2(\R^2)$ solve the one-dimensional wave equation. If $u$ and
$u_t$ vanish on the initial interval $[x_0-t_0,x_0+t_0]$, where $t_0>0$, then
\[
u(x,t)=0 \quad\text{whenever}\quad 0\leq t\leq t_0,
\qquad |x-x_0|<t_0-t.
\]
\end{lemma}
\begin{proof}
\leanok
For each $0\leq s\leq t$, the two characteristic feet $(x+s,0)$ and
$(x-s,0)$ remain in the initial interval whenever the displayed cone
inequalities hold. The local zero initial displacement at each foot implies
that the corresponding initial spatial derivative is zero; together with the
zero initial velocity, both characteristic Riemann invariants vanish there.
Propagation along the two characteristics gives $u_t(x,s)=0$ on the whole
vertical path. Integrating in time and using the zero initial value gives
$u(x,t)=0$.
\end{proof}

\begin{lemma}[Smoothness of spherical means]
\label{lem:wave-spherical-mean-smoothness}
\group{Wave Equation}
\lean{EvansLib.waveSphericalMean_contDiff}
\leanok
If $u\in C^\infty(\R^n\times\R)$ and $x\in\R^n$, then
\[
(r,t)\longmapsto
\int_{\partial B(0,1)}u(x+r\omega,t)\,dS(\omega)
\]
is smooth on $\R^2$.
\end{lemma}
\begin{proof}
\leanok
For fixed $\omega$, the integrand is smooth in $(r,t)$. Every iterated
parameter derivative is continuous jointly in $(r,t,\omega)$, and the unit
sphere is compact. Differentiation under the integral therefore applies at
every order and proves smoothness.
\end{proof}

\begin{lemma}[Time derivatives of spherical means]
\label{lem:wave-spherical-mean-time-derivative}
\group{Wave Equation}
\uses{lem:wave-spherical-mean-smoothness}
\lean{EvansLib.iteratedFDeriv_waveSphericalMean_time_two_eq_laplacian}
\leanok
Suppose $u\in C^\infty(\R^n\times\R)$ satisfies $u_{tt}=\Delta_xu$.
For every $x\in\R^n$,
\[
\frac{\partial^2}{\partial t^2}
\int_{\partial B(0,1)}u(x+r\omega,t)\,dS(\omega)
=\int_{\partial B(0,1)}\Delta_xu(x+r\omega,t)\,dS(\omega).
\]
\end{lemma}
\begin{proof}
\leanok
Differentiate twice under the integral, which is justified by
\cref{lem:wave-spherical-mean-smoothness}, and substitute
$u_{tt}=\Delta_xu$ pointwise.
\end{proof}

\begin{lemma}[Weak divergence-form radial equation]
\label{lem:weak-divergence-form-radial-equation}
\group{Harmonic Functions}
\uses{lem:sourced-weak-radial-ode-identity}
\lean{EvansLib.integral_unitSphereRadialIntegral_mul_deriv_weightedDeriv_eq_source}
\leanok
Under the hypotheses of
\cref{lem:sourced-weak-radial-ode-identity}, let
$\phi\in C_c^\infty((0,R))$. Then
\[
\int_0^\infty S_u(r)
  \frac{d}{dr}\left(r^{n-1}\phi'(r)\right)\,dr
=\int_0^\infty r^{n-1}S_{\Delta u}(r)\phi(r)\,dr.
\]
\end{lemma}
\begin{proof}
\leanok
Extend $g(\rho)=\phi(\sqrt\rho)$ smoothly across $\rho=0$; this is possible
because the support of $\phi$ is compactly contained in $(0,R)$. For $r>0$,
$g(r^2)=\phi(r)$, and two applications of the chain rule give
\[
r^{n-1}\left(4r^2g''(r^2)+2n g'(r^2)\right)
=\frac{d}{dr}\left(r^{n-1}\phi'(r)\right).
\]
Substitute these identities into
\cref{lem:sourced-weak-radial-ode-identity}.
\end{proof}

\begin{lemma}[Weak Green identity for harmonic functions]
\label{lem:weak-green-identity-compact-test}
\group{Harmonic Functions}
\lean{EvansLib.integral_mul_laplacian_eq_zero_of_harmonicOnNhd}
\leanok
Let $U \subset \R^n$ be open and let $u$ be harmonic in a neighborhood of $U$.
If $w \in C_c^\infty(\R^n)$ and the closed support of $w$ is contained in $U$, then
\[
\int_{\R^n} u\,\Delta w\,dy = 0.
\]
\end{lemma}
\begin{proof}
\leanok
Apply one-dimensional integration by parts twice in each coordinate direction.
The compact support of $w$ and its derivatives removes the boundary terms, while
the support inclusion places every derivative of $u$ that occurs in the calculation
inside $U$. Summing over the coordinate directions gives
\[
\int u\,\Delta w = \int (\Delta u)w = 0,
\]
because $u$ is harmonic on $U$.
\end{proof}

\begin{lemma}[Weak maximum principle with anisotropic heat regularity]
\label{lem:weak-maximum-principle-heat-anisotropic}
\group{Maximum Principles}
\uses{def:parabolic-cylinder-heat-ball,def:parabolic-boundary-heat-ball}
\lean{EvansLib.exists_parabolicBoundary_isMaxOn_anisotropic}
\leanok
For a continuous solution on a bounded parabolic cylinder, the maximum is
attained on the parabolic boundary under only one time derivative and two pure
spatial derivatives on each coordinate line at interior times. The heat equation
is imposed in the corresponding line-derivative form, so no mixed or second-time
derivatives are required.
\end{lemma}
\begin{proof}
\leanok
Fix $0<T'<T$ and $\varepsilon>0$, and put $v(x,t)=u(x,t)-\varepsilon t$.
Compactness of $\overline{U_{T'}}$ gives a point $q$ where $v$ is maximal. If
$q\in U_{T'}$, then every spatial coordinate section has a local maximum at
$q$, so $v_{x_jx_j}(q)\leq0$ for each $j$. The time section has a maximum from
the left, hence $v_t(q)\geq0$. Consequently
\[
 0\leq v_t(q)-\sum_jv_{x_jx_j}(q).
\]
But the line-form heat equation for $u$ gives
\[
 v_t(q)-\sum_jv_{x_jx_j}(q)=-\varepsilon<0,
\]
a contradiction. Thus the maximum of $v$ lies on $\Gamma_{T'}$.

Let $z_b$ maximize $u$ on the compact full parabolic boundary $\Gamma_T$.
Since $\Gamma_{T'}\subseteq\Gamma_T$ and all its time coordinates are
nonnegative, the preceding conclusion gives, for
$p\in\overline{U_{T'}}$,
\[
 u(p)\leq u(z_b)+\varepsilon T.
\]
If $u(p)>u(z_b)$, choosing
$\varepsilon=(u(p)-u(z_b))/(2T)$ is a contradiction; hence
$u(p)\leq u(z_b)$. Every point with time coordinate less than $T$ belongs to
some shorter closed cylinder. A point at time $T$ is approached from below by
points of $\overline{U_T}$, so continuity gives the same inequality there.
Therefore $z_b$ is a parabolic-boundary maximizer on the full closed cylinder.
\end{proof}

\begin{lemma}[Weak radial ODE identity for harmonic functions]
\label{lem:weak-radial-ode-identity}
\group{Harmonic Functions}
\uses{lem:weak-green-identity-compact-test}
\lean{EvansLib.integral_mul_squaredRadialLaplacianProfile_eq_zero_of_harmonicOnNhd}
\leanok
Assume $n\geq 1$. Let $u$ be harmonic on an open neighborhood of
$\overline{B(0,R)}$. For a
smooth one-variable profile $g$ that vanishes for $\rho>R^2$, set
$w(z)=g(|z|^2)$. Then
\[
\int_{\R^n}u(z)\left(4|z|^2g''(|z|^2)+2n g'(|z|^2)\right)\,dz=0.
\]
This is the centered smooth-test identity preceding the sharp ball mean-value
formula.
\end{lemma}
\begin{proof}
\leanok
The squared-radial test is compactly supported in $\overline{B(0,R)}$.
Away from the origin its Laplacian is
$4|z|^2g''(|z|^2)+2n g'(|z|^2)$; the origin is a null set in positive
dimension. Apply the weak Green identity to this test function.
\end{proof}

\begin{theorem}[Converse representation of a one-dimensional wave]
\label{thm:wave-1d-converse-representation}
\group{Wave Equation}
\uses{thm:dalembert-initial-conditions,lem:wave-cauchy-uniqueness-1d}
\lean{EvansLib.wave_exists_general_representation_1d}
\leanok
Conversely, every global $C^2$ solution of the one-dimensional wave equation can
be written
$$u(x,t) = F(x+t) + G(x-t)$$
for some $F,G \in C^2(\R)$.
\end{theorem}
\begin{proof}
\leanok
Take the initial displacement $g(x)=u(x,0)$ and initial velocity $h(x)=u_t(x,0)$.
Their regularity is inherited from $u$. The d'Alembert solution with data $(g,h)$
has the same initial displacement and velocity, so one-dimensional Cauchy
uniqueness identifies it with $u$. If $H(y)=\int_0^y h(s)\,ds$, then $H\in C^2$
and d'Alembert's formula becomes
\[
u(x,t)=\frac{g(x+t)+H(x+t)}2+\frac{g(x-t)-H(x-t)}2.
\]
Thus $F=(g+H)/2$ and $G=(g-H)/2$ are $C^2$ profiles with the required form.
\end{proof}
